% ============================================================================
% Fichier : annexes/Ann_E_bibliographie.tex
% Titre   : Bibliographie Commentée
% Version : 1.0 -- novembre 2026
% ============================================================================
% Uniquement des entrées de bibliographie/GCI_references.bib.
% ============================================================================

\chapter{Bibliographie Commentée}
\label{ann:biblio}

Pour chaque partie, quelques lectures, avec ce qu'elles apportent. Les
références complètes figurent dans la bibliographie du livre. Les
publications d'ISACA et d'AXELOS sont paraphrasées dans le cours~; leur
lecture directe reste indispensable.

\section*{Prologue et Partie I · Fondements}

\begin{itemize}
    \item Weill et Ross~\cite{weill2004itgovernance}~: les droits de décision
          et les archétypes de gouvernance~; le livre fondateur.
    \item Jensen et Meckling~\cite{jensen1976agency}, Freeman~\cite{freeman1984stakeholder},
          Davis et al.~\cite{davis1997stewardship}~: agence, parties
          prenantes, intendance~; les théories sous la gouvernance.
    \item Cadbury~\cite{cadbury1992}, principes de l'OCDE~\cite{oecd2023principles},
          King~IV~\cite{iod2016king4}~: les codes de gouvernance d'entreprise.
    \item Henderson et Venkatraman~\cite{henderson1993alignment}~: le modèle
          d'alignement stratégique.
    \item Les sources de première main sur les défaillances~: SEC sur Knight
          Capital~\cite{sec2013knight}, FCA et PRA sur TSB~\cite{fca2022tsb},
          analyses de l'incident CrowdStrike~\cite{crowdstrike2024rca,microsoft2024crowdstrike}.
    \item Thorp~\cite{thorp1998paradox} et Weill et Broadbent~\cite{weill1998infrastructure}~:
          valeur et portefeuille.
\end{itemize}

\section*{Partie II · COBIT 2019}

\begin{itemize}
    \item ISACA, \emph{Introduction and Methodology}~\cite{isaca2018cobitintro}~:
          principes, composantes, facteurs de conception, capacité.
    \item ISACA, \emph{Governance and Management Objectives}~\cite{isaca2018cobitobjectifs}~:
          les quarante objectifs, à consulter objectif par objectif.
    \item COBIT~4.1~\cite{itgi2007cobit41}~: encore cité dans beaucoup de
          dossiers~; les modèles de maturité.
    \item De Haes et Van Grembergen~\cite{dehaes2015egit}~: la gouvernance
          des SI d'entreprise, entre recherche et pratique.
\end{itemize}

\section*{Partie III · Cadre normatif, juridique et déontologique}

\begin{itemize}
    \item ISO/IEC~38500~\cite{iso38500,iso38500_2024} et
          ISO/IEC~27000~\cite{iso27000_2018}.
    \item COSO ERM~\cite{coso2004ermfr,coso2017erm}~: la gestion des risques
          de l'entreprise.
    \item IIA~: CRIPP~\cite{iia2017cripp}, Normes mondiales~\cite{iia2024gias},
          modèle des trois lignes~\cite{iia2020troislignes}~; INTOSAI,
          GUID~5100~\cite{intosai2019guid5100}.
    \item Textes~: loi n°~2010/012~\cite{loi2010012}, COBAC
          R-2016/04~\cite{cobac2016r04}, AUSCGIE~\cite{ohada2014auscgie},
          régime financier de l'État~\cite{minfi2023regimefinancier},
          Sarbanes-Oxley~\cite{sox2002}~; voir l'Annexe~\ref{ann:textes}.
\end{itemize}

\section*{Partie IV · ITIL 4 et pratique opérationnelle}

\begin{itemize}
    \item AXELOS, ITIL~4~\cite{axelos2019itil4}~: le système de valeur des
          services et les pratiques.
    \item PeopleCert, ITIL (2026)~\cite{peoplecert2026itil5}~: l'évolution
          du référentiel, à suivre.
\end{itemize}

\section*{Partie V · Défis avancés}

\begin{itemize}
    \item Manifeste agile~\cite{agile2001manifeste} et Forsgren et
          al.~\cite{forsgren2018accelerate}~: agilité et performance de la
          livraison.
    \item NIST, définition du nuage~\cite{nist2011cloud}.
    \item NIST AI RMF~\cite{nist2023airmf}, Recommandation de
          l'UNESCO~\cite{unesco2021ethique}, règlement européen sur
          l'IA~\cite{ue2024aiact}~: gouverner l'IA.
\end{itemize}

\section*{Partie VI · Recherche}

\begin{itemize}
    \item Sambamurthy et Zmud~\cite{sambamurthy1999arrangements}, Peterson~\cite{peterson2004crafting},
          De Haes et Van Grembergen~\cite{dehaes2009exploratory}, Brown et
          Grant~\cite{brown2005framing}, Wu et al.~\cite{wu2015mechanisms}~:
          les arrangements et mécanismes de gouvernance.
    \item Becker et al.~\cite{becker2009maturity}, Pöppelbuß et
          Röglinger~\cite{poppelbuss2011maturity}, Simonsson et
          al.~\cite{simonsson2010maturity}~: concevoir et éprouver un modèle
          de maturité.
    \item Hevner et al.~\cite{hevner2004design}~: la science du design en
          systèmes d'information.
\end{itemize}
